\section{Comparative Analysis of Contemporary Legal AI Systems}
\label{sec:comparison}

\subsection{A common framework for comparing legal AI architectures}

Contemporary legal AI systems differ not only in the models they use but also in how they retrieve,
organize, interpret, and verify legal evidence. A fair comparison must therefore examine the complete
processing pipeline rather than focusing exclusively on the underlying language model or a reported
accuracy score. This section compares representative systems using six dimensions:
\begin{enumerate}[leftmargin=1.4em, itemsep=2pt, topsep=2pt]
\item \textbf{Architecture:} how the system organizes its processing components.
\item \textbf{Retrieval strategy:} how it identifies and ranks relevant legal material.
\item \textbf{Generation method:} how it produces classifications, legal analysis, or drafted text.
\item \textbf{Evidence verification:} whether it checks source validity and claim-level support.
\item \textbf{Evaluation:} how its performance is measured.
\item \textbf{Limitations:} what its architecture or experimental design does not establish.
\end{enumerate}

The comparison distinguishes between legal language models, benchmarks, retrieval systems, and
end-to-end applications. This distinction is necessary because a benchmark such as LegalBench-RAG
evaluates a specific capability, whereas a platform such as JusBuild is designed to support a legal
workflow. Their contributions are complementary, but they cannot be ranked directly without common
experiments. A benchmark for legal classification cannot establish that a system is reliable at
retrieving current statutes, and a document-drafting system cannot automatically be considered
superior to a legal search system.

\subsection{Architecture and retrieval strategy}

Table~\ref{tab:arch8} compares the architecture and retrieval strategy of representative legal AI
approaches.

\begin{table*}[tp]
\centering
\caption{Architectural comparison of representative legal AI approaches.}
\label{tab:arch8}
\footnotesize
\setlength{\tabcolsep}{5pt}
\renewcommand{\arraystretch}{1.3}
\begin{tabularx}{\textwidth}{>{\raggedright\arraybackslash}p{2.8cm} >{\raggedright\arraybackslash}X >{\raggedright\arraybackslash}X >{\raggedright\arraybackslash}X}
\toprule
\textbf{System} & \textbf{Architecture} & \textbf{Retrieval strategy} & \textbf{Generation or output method} \\
\midrule
LEGAL-BERT \citep{chalkidis2020} & Legal-domain pretrained transformer & Not inherently a retrieval system & Contextual representations for downstream legal NLP tasks \\
LexGLUE \citep{chalkidis2022} & Benchmark suite for legal language understanding & Depends on the evaluated model and task & Task-specific predictions, such as classification \\
LegalBench \citep{guha2023} & Multi-task legal reasoning benchmark & Not a single retrieval architecture & Model-generated answers or task outputs evaluated against task-specific criteria \\
LegalBench-RAG \citep{pipitone2024} & Retrieval benchmark with annotated legal text segments & Evaluates retrieval of relevant legal passages & Evaluates retrieval rather than prescribing one answer-generation architecture \\
JusBuild \citep{buffa2025} & Structured legal document-building workflow & Document segmentation and vector-based retrieval & Combines retrieved precedent material with AI-assisted drafting and human oversight \\
Pakistani legal RAG framework \citep{ref312025} & Jurisdiction-specific question-answering pipeline & Hybrid keyword and semantic retrieval & LLM-generated legal answers grounded in retrieved materials \\
Ren et al.\ \citep{ren2026} & Legal LLM improvement pipeline incorporating data synthesis and RAG & Retrieval-augmented generation alongside synthetic-data development & Legal LLM outputs assessed through benchmark and model-evaluation procedures \\
Legalyze \citep{hamza2026} & Structure-aware, agentic legal RAG platform & Hierarchical indexing, lexical and dense retrieval, reciprocal-rank fusion & Agentic legal analysis and synthesis grounded in retrieved legal evidence \\
\bottomrule
\end{tabularx}
\end{table*}

The table demonstrates that architectural differences correspond to different research objectives.
Domain-adapted transformers aim to improve the model's understanding of legal language. Retrieval
benchmarks measure whether relevant evidence can be found. End-to-end RAG applications combine
retrieval with generation, while structure-aware systems additionally attempt to preserve
relationships within legal documents.

These differences matter because a model that performs well on a legal classification task may not
retrieve the relevant statutory provision, and a system that retrieves the right passage may still
misinterpret its legal meaning. The architecture must therefore be assessed against the particular
failure modes it is intended to address.

\subsection{Evidence verification and evaluation}

Table~\ref{tab:evid8} summarizes evidence handling and evaluation considerations across the same
systems.

\begin{table*}[tp]
\centering
\caption{Evidence handling and evaluation considerations.}
\label{tab:evid8}
\footnotesize
\setlength{\tabcolsep}{5pt}
\renewcommand{\arraystretch}{1.3}
\begin{tabularx}{\textwidth}{>{\raggedright\arraybackslash}p{2.8cm} >{\raggedright\arraybackslash}X >{\raggedright\arraybackslash}X >{\raggedright\arraybackslash}X}
\toprule
\textbf{System or approach} & \textbf{Evidence verification} & \textbf{Evaluation focus} & \textbf{Key limitation to consider} \\
\midrule
LEGAL-BERT & No inherent citation-verification mechanism & Performance on downstream legal NLP tasks & Legal language competence does not establish evidence-grounded reasoning \\
LexGLUE & Determined by the evaluated task and model & Standardized legal language understanding & Task scores do not directly measure citation support \\
LegalBench & Task-dependent; not a universal citation verifier & A broad range of legal reasoning tasks & Performance may vary considerably by task and legal concept \\
LegalBench-RAG & Evaluates retrieval against annotated relevant passages & Passage-level retrieval quality & Does not, by itself, establish the correctness of a generated legal answer \\
JusBuild & Uses a structured drafting workflow with human oversight & Legal document-building performance in its evaluated settings & Human review and drafting performance do not establish general-purpose legal reliability \\
Pakistani legal RAG framework & Evidence is retrieved for grounded answering; verification should be assessed separately & Legal question answering over selected Pakistani legal materials & Grounding alone does not establish source completeness, temporal validity, or claim-level entailment \\
Ren et al. & Retrieval and evaluation components support the broader legal LLM pipeline & Legal model performance and evaluation of generated outputs & Results depend on the selected benchmarks and experimental settings \\
Legalyze & Citation-grounded response design and a hybrid retrieval pipeline & Retrieval and answer reliability within its reported evaluation setting & Performance on one corpus and query set cannot establish universal legal correctness \\
\bottomrule
\end{tabularx}
\end{table*}

A critical distinction emerges: citation presence, citation validity, and citation support are three
different properties. A system may generate a citation, point to an existing legal source, and still
misrepresent what that source establishes. Accordingly, evaluation should measure whether citations
resolve correctly and whether the cited passages support individual claims.

The same caution applies to numerical results. Retrieval recall, legal answer accuracy,
unsupported-claim rate, and citation precision measure different properties. Scores reported on
different datasets or under different annotation protocols should not be compared as if they were
measurements of the same outcome. For this reason, contemporary systems should ideally be evaluated
through a common protocol combining passage-level retrieval metrics, expert-assessed legal
correctness, citation-support analysis, temporal and jurisdictional checks, and operational measures
such as latency and cost.

\subsection{Comparing system design philosophies}

The systems above illustrate three broad design philosophies.

\subsubsection{Model-centric legal AI}

Model-centric approaches improve legal performance by adapting language models to legal text, tasks, or
reasoning patterns. LEGAL-BERT and research evaluated using LexGLUE exemplify this direction. LegalBench
complements it by testing a wider range of legal reasoning capabilities.

The principal advantage is that domain adaptation can improve a model's ability to interpret legal
terminology, identify relevant concepts, and perform specialized NLP tasks. However, model-centric
systems remain limited when a question depends on an authoritative source that was absent from training,
has changed since training, or must be cited precisely.

The key distinction is between understanding legal language and retrieving legally applicable evidence.
A model can recognize the meaning of a statutory provision without establishing whether that provision
is current or whether it applies to the user's specific question.

\subsubsection{Retrieval-centric legal AI}

Retrieval-centric approaches place greater emphasis on locating relevant legal material. LegalBench-RAG
demonstrates the importance of retrieving precise text segments rather than returning only broad
document-level matches. The Pakistani legal RAG framework illustrates the use of hybrid retrieval to
combine semantic similarity with keyword matching.

The main advantage is that answers can be conditioned on external legal sources, reducing dependence on
the model's internal knowledge. Retrieval can also make responses more auditable because a user can
inspect the source passages.

Nevertheless, retrieval quality depends on corpus coverage, document parsing, chunking, query
interpretation, and ranking. Even a highly relevant passage can be misinterpreted or applied outside its
proper legal context. Retrieval-centric systems therefore need a separate mechanism for validating the
relationship between evidence and generated claims.

\subsubsection{Structure-aware and workflow-oriented legal AI}

A third direction combines retrieval with legal document structure, task-specific processing, or
coordinated workflows. JusBuild uses document segmentation and section-level retrieval to assist with
structured drafting. Legalyze combines hierarchical legal indexing with lexical and dense retrieval,
rank fusion, and agentic analysis.

These approaches recognize that legal documents are not simply collections of independent text passages.
A provision belongs to a section, chapter, statute, or legal instrument; a judgment contains procedural
history, arguments, reasoning, and a decision. Preserving this organization can help a system retrieve
the correct passage while retaining the context needed to interpret it.

Workflow-oriented designs can additionally separate retrieval, analysis, synthesis, and verification.
Their potential benefit is improved traceability and more explicit control over intermediate operations.
Their cost is architectural complexity: every additional stage introduces implementation, latency, and
evaluation requirements. The critical question is not whether a system has more components, but whether
those components address identifiable failure modes and produce measurable improvements.

\subsection{Cross-system comparison: capabilities and unresolved gaps}

The comparison can be extended beyond individual systems to the capabilities that a dependable legal AI
platform should support, as set out in Table~\ref{tab:capabilities}.

\begin{table*}[tp]
\centering
\caption{Capabilities of a dependable legal AI platform and associated research directions.}
\label{tab:capabilities}
\footnotesize
\setlength{\tabcolsep}{5pt}
\renewcommand{\arraystretch}{1.3}
\begin{tabularx}{\textwidth}{>{\raggedright\arraybackslash}p{3.0cm} >{\raggedright\arraybackslash}X >{\raggedright\arraybackslash}X >{\raggedright\arraybackslash}X}
\toprule
\textbf{Capability} & \textbf{Why it matters} & \textbf{Typical weakness if omitted} & \textbf{Research direction} \\
\midrule
Legal language understanding & Interprets legal terminology and complex syntax. & Relevant legal meaning may be misunderstood. & Legal-domain adaptation and task-specific evaluation. \\
Hybrid retrieval & Finds both exact legal expressions and semantically related passages. & Statutory identifiers or paraphrased queries may be missed. & Combine lexical and dense retrieval. \\
Rank fusion and reranking & Improves the ordering of candidate evidence. & Relevant passages may remain buried in the result list. & Evaluate fusion and reranking through ablations. \\
Hierarchical indexing & Preserves relationships between passages and their parent documents. & Retrieved fragments may lose statutory or judgment context. & Structure-aware segmentation and metadata. \\
Knowledge-graph retrieval & Represents explicit relations between legal entities and authorities. & Cross-document relationships may remain difficult to retrieve. & Graph construction, entity linking, and graph--text retrieval. \\
Citation verification & Tests whether cited sources exist and support generated claims. & Responses may contain invalid or misleading citations. & Claim-level evidence verification. \\
Temporal and jurisdictional control & Helps ensure that retrieved law is applicable. & Superseded or foreign legal authorities may be used. & Versioned corpora and jurisdiction-aware filtering. \\
Multilingual retrieval & Supports queries expressed in different languages or scripts. & Users may fail to retrieve evidence phrased differently from their query. & Cross-lingual retrieval, transliteration, and legal terminology alignment. \\
Reproducible evaluation & Makes performance claims independently testable. & Improvements may reflect dataset or evaluation differences. & Public test sets, transparent protocols, and statistical analysis. \\
\bottomrule
\end{tabularx}
\end{table*}

This comparison reveals that the main challenge is not necessarily the absence of an individual
technique. Rather, the challenge is integrating these techniques into a coherent pipeline while
preserving evidence provenance and enabling meaningful evaluation. For example, hybrid retrieval and
reranking can improve candidate selection, but neither guarantees that the selected passage is legally
authoritative. Hierarchical indexing can preserve document context, but it does not automatically
establish relationships between separate judgments and statutes. A knowledge graph can encode those
relationships, but its usefulness depends on the accuracy, completeness, and provenance of the graph.
Citation verification can identify unsupported claims, but only if the verification process has access
to suitable evidence and appropriate legal criteria.

\subsection{The Pakistani legal AI context}

The Pakistani legal context provides a useful case study because jurisdiction-specific legal materials
must be interpreted within their own institutional and documentary structures. A system developed
primarily around US or European legal resources cannot be assumed to perform reliably on Pakistani
statutes and judgments without appropriate corpus construction, jurisdictional adaptation, and
evaluation.

\citet{ref312025} describe a RAG-based framework for Pakistani legal question answering that combines
semantic and keyword-based retrieval over selected legal materials, including the Code of Criminal
Procedure and Pakistan Penal Code. This work illustrates the importance of combining legal-domain
evidence with retrieval methods suited to the source material.

Legalyze \citep{hamza2026} follows a related but more structure-aware approach. Its architecture
incorporates hierarchical indexing, lexical and dense retrieval, reciprocal-rank fusion, and agentic
legal analysis over Pakistani legal materials. These design choices address complementary challenges:
lexical retrieval supports exact legal terminology, dense retrieval supports semantic variation, rank
fusion combines result lists, and hierarchical indexing helps preserve legal document organization.

However, a meaningful comparison between the two systems would require a shared evaluation dataset,
consistent query sets, documented corpus coverage, equivalent retrieval budgets, and common measures of
retrieval and answer quality. The existence of hybrid retrieval in both systems is not sufficient to
establish which performs better. Likewise, differences in reported hallucination rates should not be
attributed to a single architectural component unless the experimental design isolates that component.

A further research opportunity concerns multilingual access. Pakistani legal materials and user
questions may involve English, Urdu, and Roman Urdu, while statutory text and case-law terminology may
not map neatly across languages. A robust system must therefore align the user's intended legal concept
with the authoritative terminology used in the source corpus, rather than relying exclusively on literal
translation.

\subsection{A unified framework for comparing legal AI systems}

The preceding analysis suggests that future comparisons should separate four levels of system
performance:
\begin{enumerate}[leftmargin=1.4em, itemsep=3pt, topsep=2pt]
\item \textbf{Retrieval performance.} Does the system find the relevant legal material and rank it
appropriately? \textit{Measures:} Recall@$k$, MRR, nDCG, passage-level precision.
\item \textbf{Evidence integrity.} Does the system preserve source identity, document context, temporal
validity, and jurisdiction? \textit{Measures:} source resolution, provenance completeness, temporal and
jurisdictional error rates.
\item \textbf{Legal answer quality.} Does the answer correctly interpret and apply the retrieved
authority? \textit{Measures:} expert-assessed correctness, claim support, citation accuracy,
unsupported-claim rate.
\item \textbf{Operational performance.} Can the system provide useful results within acceptable
computational and review costs? \textit{Measures:} latency, computational cost, failure recovery, and
human review burden.
\end{enumerate}

This framework discourages the use of one aggregate score to conceal important weaknesses. A system may
retrieve highly relevant passages but generate inaccurate legal conclusions; another may generate
convincing answers but rely on incomplete or outdated evidence. Reporting each level separately makes
such trade-offs visible. Future comparative studies should also report whether the system is
open-source, whether its corpus and evaluation data are accessible, whether its outputs can be
reproduced, and whether its architecture supports independent component replacement.

\subsection{What the comparison reveals}

Three conclusions follow from the comparison. First, no single architectural family addresses every
legal AI problem. Domain-specific language models improve legal language processing, retrieval systems
improve access to external evidence, and structured workflows can coordinate more specialized
operations.

Second, evidence verification remains distinct from evidence retrieval. Lexical retrieval, dense
retrieval, reranking, and rank fusion can improve the candidate evidence presented to a model. They do
not automatically establish that the final answer is supported by that evidence.

Third, jurisdiction and document structure should be treated as architectural requirements rather than
optional metadata. A system that retrieves semantically similar passages without considering the
governing jurisdiction, the document hierarchy, or the effective date of a provision can return
plausible but legally inapplicable evidence. These observations motivate the next section, which
examines jurisdiction-specific reasoning, legal hierarchy, and multilingual access.
